% % seminario-A8.tex
% Slides
\begin{flushleft}
    \vspace{-1.75em}
    {\color{gray} Presentado por: Eric Lucero Gonzalez \hfill eric.lucero@alumnos.upm.es} \\
    {\color{gray} Impartido por: Angélica de Antonio \\Aula: A-6205}
    % \section*{Seminario A8: Protección de resultados científicos y Transferencia Tecnológica}
\end{flushleft}
\vspace{-4em}
% --- Actividad I ---
% \begin{boxA}{Activity I: Results and impact of your research}
%     \textbf{Instructions:}\\
%     For your main research activity, identify:
%     \begin{enumerate}
%         \item \textbf{Needs} to be addressed.
%         \item \textbf{Expected output(s)} (knowledge, data, product, process, service, activity).
%         \item \textbf{Expected outcome(s)}.
%         \item \textbf{Who can benefit}.
%         \item \textbf{Expected impact:} Scientific, Economic, and Social.
%     \end{enumerate}
% \end{boxA}

% --- Actividad I ---
% \begin{boxB}{Activity I Solution: Results and impact of your research}
%     For your main research activity, identify:
%     \begin{enumerate}
%         \item \textbf{Needs} to be addressed:\\
%               The shift scheduling and task assignation to the onboard line-handlers workers in the Panama Canal.
%               Manual or semi-automated scheduling of vessel transits through lock chambers is exceptionally complex due to high traffic volume and stringent operational constraints.
%               There is a critical need to maximize vessel throughput, minimize waiting times, and eliminate human error in the logistics planning of the Pacific locks.
%         \item \textbf{Expected output(s):}
%               \begin{itemize}
%                   \item \textbf{Knowledge:} A formal mathematical model of the specific maritime timetabling problem and the development of a novel metaheuristic.
%                   \item \textbf{Product:} A Decision Support System (DSS) software tool designed for direct implementation by the Panama Canal Authority (ACP).
%               \end{itemize}
%         \item \textbf{Expected outcome(s):}
%               \begin{itemize}
%                   \item Optimization of existing infrastructure utilization without requiring physical expansion.
%                   \item Systematic reduction in average vessel Transit Time (TT).
%                   \item Enhanced resilience through dynamic re-scheduling capabilities to respond to unforeseen operational disruptions.
%               \end{itemize}
%         \item \textbf{Who can benefit:}
%               \begin{itemize}
%                   \item \textbf{Panama Canal Authority (ACP):} Improved operational efficiency and strategic resource management.
%                   \item \textbf{Shipping Companies:} Significant cost savings derived from reduced idle time and fuel consumption.
%                   \item \textbf{Scientific Community:} New contributions to the fields of operations research and computational intelligence, providing a scalable metaheuristic applicable to other port logistics challenges.
%               \end{itemize}
%         \item \textbf{Expected impact:}
%               \begin{itemize}
%                   \item \textbf{Scientific:} Advancing the state-of-the-art in evolutionary computing applied to Large-scale Timetabling Problems (LSTP).
%                   \item \textbf{Economic:} Increased profitability for the Canal by enabling higher daily vessel capacity and reducing operational overhead.
%                   \item \textbf{Social:} Strengthening the global supply chain (a vital maritime node) and contributing to environmental sustainability by reducing the carbon footprint of waiting vessels.
%               \end{itemize}
%     \end{enumerate}
% \end{boxB}
% \clearpage
% --- Actividad II ---
% \begin{boxA}{Activity II: Communication of results}
%     \textbf{Instructions:}\\
%     Create two Communication Actions (2 slides) according to these guidelines:
%     \begin{itemize}
%         \item \textbf{Rationale (WHY):} Create awareness about your work.
%         \item \textbf{Audiences (WHO):} Classmates (technical) and Family/Friends.
%         \item \textbf{Messages (WHAT):} Outcomes, results, and timeline.
%         \item \textbf{Channels (HOW):} Format for one slide.
%         \item \textbf{Timeline (WHEN):} Current status.
%     \end{itemize}
% \end{boxA}
% --- Actividad II ---
% \begin{boxB}{Activity II - Communication of results (Classmates)}
%     \begin{itemize}
%         \item \textbf{Rationale (WHY):}\\
%               Operational efficiency in the Panama Canal is often hindered by labor-intensive scheduling bottlenecks.
%               A critical component is the management of line-handlers\footnote{Personnel responsible for securing vessels within the locks.}.
%               This research addresses the multi-objective Shift Scheduling and Rostering Problem (SSRP) for line-handlers, aiming to optimize task assignment while simultaneously satisfying complex constraints such as worker shift preferences, mandatory rest periods, and leave requests.
%         \item \textbf{Audiences (WHO):}\textbf{Peer Researchers and Classmates:}
%               %   Focus on the mathematical modeling of constraints and the algorithmic architecture of the proposed metaheuristic.
%         \item \textbf{Messages (WHAT):}
%               \begin{itemize}
%                   \item \textbf{Expected outcomes:} An automated framework that significantly reduces scheduling lead time, eliminates constraint violations, and achieves a Pareto-optimal balance between operational needs and staff well-being.
%                   \item \textbf{Current results:} Formalization of the constraint satisfaction model and preliminary benchmarking of the metaheuristic core.
%                   \item \textbf{Timeline:} Prototype finalization and validation using historical ACP operational data within 12 months.
%               \end{itemize}
%         \item \textbf{Channels (HOW):} Technical presentation featuring algorithmic flowcharts and comparative performance metrics (Manual vs. Algorithmic).
%         \item \textbf{Timeline (WHEN):} Immediate / Seminar session.
%     \end{itemize}
% \end{boxB}

% \begin{boxB}{Activity II - Communication of results (Family and Friends)}
%     \begin{itemize}
%         \item \textbf{Rationale (WHY):}\\
%               Imagine the chaos of trying to organize a 24/7 work schedule for over a thousand people, all while following strict safety rules
%               (like mandatory 12-hour breaks) and dealing with a fluctuating number of massive ships arriving every day. Currently, this ``giant puzzle''
%               is solved largely by hand, which is slow and prone to mistakes.
%         \item \textbf{Audiences (WHO):} \textbf{Family and Friends:}
%               %   Focus on the real-world impact—how smart software makes life easier for workers and helps the Canal run smoothly.

%         \item \textbf{Messages (WHAT):}
%               \begin{itemize}
%                   \item \textbf{Expected outcomes:} I am building an ``intelligent assistant'' that uses advanced math to find the best possible schedule in seconds, making sure workers are less stressed and the Canal stays efficient.
%                   \item \textbf{Current results:} I have defined all the ``rules of the game'' and started testing the first version of the brain behind the software.
%                   \item \textbf{Timeline:} The final tool will be ready to be tested with real Canal data in about a year.
%               \end{itemize}
%         \item \textbf{Channels (HOW):} Casual conversation using analogies (like solving a Rubik's Cube or a complex puzzle) and simple visual comparisons.
%         \item \textbf{Timeline (WHEN):} Next social gathering / Now.
%     \end{itemize}
% \end{boxB}
% \clearpage

% --- Actividad III ---
% \begin{boxA}{Activity III: Exploitation of results}
%     \textbf{Instructions:}\\
%     Analyze your research's commercial path:
%     \begin{itemize}
%         \item Identify \textbf{Potential Users / Clients}.
%         \item Do results need to be protected?
%         \item Analyze applicability of different Exploitation Strategies:
%               \begin{itemize}
%                   \item Own Exploitation
%                   \item Licensing to third parties - Royalties
%                   \item Creation of spin-offs
%               \end{itemize}
%     \end{itemize}
% \end{boxA}
% % --- Actividad III ---
% \begin{boxB}{Activity III: Exploitation of results}
%     \textbf{Instructions:}\\
%     For your main research activity:
%     \begin{itemize}
%         \item \textbf{Identify Potential Users / Clients:}
%               \begin{itemize}
%                   \item The \textbf{Panama Canal Authority (ACP)} as the primary client and final user.
%                   \item Universidad Politécnica de Madrid as the owner of the software to distribute with other PhD projects.
%                   \item International port authorities with similar scheduling management complexities.
%               \end{itemize}
%         \item \textbf{Do results need to be protected?}
%               Yes. The source code and the specific metaheuristic architecture should be protected via \textbf{Copyright} and evaluated for potential \textbf{Software Patenting} or Trade Secret.
%         \item \textbf{Analyze applicability of different Exploitation Strategies:}
%               \begin{itemize}
%                   \item \textbf{Own Exploitation:} Direct implementation within the ACP's decision support systems.
%                   \item \textbf{Licensing:} Licensing the optimization software to other maritime logistics hubs via Royalties.
%                   \item \textbf{Spin-off:} Creation of a technology-based company specialized in critical maritime infrastructure optimization.
%               \end{itemize}
%     \end{itemize}
% \end{boxB}
% \clearpage
% % --- Actividad IV ---
% \begin{boxA}{Activity IV - Maturation process}
%     \textbf{Instructions:}\\
%     Identify a technology from your doctoral project and describe its maturation process through the \textbf{TRL} (Technology Readiness Levels) scale.
% \end{boxA}

% \vspace{-1.75em}
\begin{boxB}{Activity IV - Maturation process (Solution)}
    \begin{itemize}
        \item Identify a technology from your doctoral project and describe its maturation process through the \textbf{TRL} (Technology Readiness Levels) scale.
        \item Describe its (expected) maturation process through the TRLs
    \end{itemize}
\end{boxB}

% \SmartBox{Technology:}{ An optimization algorithm for rapid optimal value finding on the timetabling line-handler shift scheduling problem.}

% Low TRLs (Research) - Gray/Default
\SmartBox{TRL 1\\ Basic Principles}{
    \textbf{Mathematical Foundation:}
    \begin{itemize}
        \item Literature review of combinatorial optimization and metaheuristic algorithms.
        \item Identification of specific constraints in the line-handler rostering problem (legal, physical, and union constraints).
    \end{itemize}
}

\vspace{-1.75em}
\SmartBox{TRL 2\\ Concept Formulation}{
    \textbf{Model Definition:}
    \begin{itemize}
        \item Mathematical formulation of the objective function (minimizing overtime/idle time).
        \item Selection of the specific metaheuristic approach (e.g., Genetic Algorithm or Simulated Annealing).
        \item Pseudocode design.
    \end{itemize}
}

\vspace{-1.75em}
\SmartBox{TRL 3\\ Proof of Concept}{
    \textbf{Analytical Lab Proof:}
    \begin{itemize}
        \item Prototype implementation in Python (Jupyter Notebooks).
        \item Verification using small, synthetic datasets to prove the algorithm converges.
        \item ``Toy problem'' solved successfully.
    \end{itemize}
}

\vspace{-1.75em}
% Mid TRLs (Development) - Green indicating growth
\SmartBox{TRL 4\\ Lab Validation}{
    \textbf{Component Validation:}
    \begin{itemize}
        \item Algorithm is encapsulated into a modular function.
        \item Tested against historical ``clean'' data.
        \item Benchmarking: Comparison of execution speed and result quality against current manual methods or baseline heuristics.
    \end{itemize}
}
\vspace{-1.75em}

\SmartBox[colorD]{TRL 5\\ Relevant Environment}{
    \textbf{Integrated Prototype:}
    \begin{itemize}
        \item Development of a basic Graphical User Interface (GUI) for non-coders.
        \item Handling of \textit{dirty} real-world data (missing fields, inconsistencies).
        \item Validation by end-users (schedulers) in a workshop setting using past scenarios.
    \end{itemize}
}

\vspace{-1.75em}
\SmartBox[colorD]{TRL 6\\ System Demonstration}{
    \textbf{\textit{Shadow Mode} Operation:}
    \begin{itemize}
        \item The system runs in parallel with the current manual process for one month.
        \item Outputs are generated but not used for operations yet.
        \item Comparison of the ``Shadow'' schedule vs. the Actual schedule to quantify efficiency gains.
    \end{itemize}
}

\vspace{-1.75em}
% High TRLs (Deployment) - Orange/Blue for Launch
\SmartBox[GithubGreen]{TRL 7\\ Operational Pilot}{
    \textbf{Field Deployment:}
    \begin{itemize}
        \item Installation on Panama Canal Authority (ACP) servers.
        \item Used for actual scheduling in a limited scope (e.g., one specific zone or shift).
        \item Real-time adjustments based on scheduler feedback during operations.
    \end{itemize}
}
\vspace{-1.75em}

\SmartBox[GithubGreen]{TRL 8\\ System Qualified}{
    \textbf{Full Commissioning:}
    \begin{itemize}
        \item Documentation, user manuals, and training completed.
        \item Final bug fixes regarding security and database integration.
        \item System is certified as the primary tool for scheduling.
    \end{itemize}
}
\vspace{-1.75em}

\SmartBox[GithubGreen]{TRL 9\\ Operations}{
    \textbf{Full Maintenance:}
    \begin{itemize}
        \item Routine daily use across all departments.
        \item Maintenance contracts and software updates active.
        \item Continuous improvement based on long-term metrics.
    \end{itemize}
}



% % --- Actividad V ---
% \begin{boxA}{Activity V: Technology attributes}
%     \textbf{Instructions:}\\
%     Analyze how these attributes act as \textbf{Facilitators} or \textbf{Barriers}:
%     \begin{itemize}
%         \item Relative Advantage
%         \item Compatibility
%         \item Complexity
%         \item Observability
%         \item Tryability.
%     \end{itemize}
%     Analyze the type of change expected for the adoption of the technology.
% \end{boxA}


% \begin{boxB}{Activity V - Technology attributes (Solution)}
%     \textbf{1. Attribute Analysis (Facilitators vs. Barriers)}
%     \begin{itemize}
%         \item \textbf{Relative Advantage (Facilitator):}
%               The algorithm reduces scheduling time from hours (manual) to minutes (automated) and creates fairer shifts. The benefit is immediate and quantifiable.

%         \item \textbf{Compatibility (Neutral/Barrier):}
%               Must respect deep-seated union rules and seniority logic. If the AI suggests a schedule that violates ``unwritten rules'' of the workforce, compatibility is low. It must be tuned to match organizational culture.

%         \item \textbf{Complexity (Barrier):}
%               Moving from Excel/Paper to a software interface may be intimidating for senior staff. The \textit{Black Box} nature of optimization (not knowing \textit{why} the AI made a decision) increases perceived complexity.

%         \item \textbf{Observability (Facilitator):}
%               Results are highly visible. A printed schedule with fewer gaps and less overtime is physically observable by management and workers immediately.

%         \item \textbf{Tryability (Facilitator):}
%               High tryability. The system can be run in \textit{Shadow Mode} (see TRL 6) without risk. Schedulers can test it on past weeks to see ``what if'' scenarios without disrupting operations.
%     \end{itemize}

%     \vspace{0.5cm}
%     \textbf{2. Type of Change Expected}
%     \begin{itemize}
%         \item \textbf{Change Magnitude:} \textit{Incremental but High-Impact.} It does not change \textit{what} they do (assigning shifts), but radically changes \textit{how} they do it (tooling).
%         \item \textbf{Level of Required Learning:} \textit{Medium.} Users need to learn the interface, but more importantly, they need to learn to ``trust'' the algorithm's suggestions.
%         \item \textbf{Time of Adjustments:} \textit{Short-term disruption.} Calibration of the algorithm takes time, but once set, the workflow becomes significantly faster.
%     \end{itemize}
% \end{boxB}
% % % --- Actividad VI ---
% % \begin{boxA}{Activity VI (Optional): Adopter analysis}
% %     \textbf{Instructions:}\\
% %     Analyze your profile according to the \textbf{Rogers model} based on:
% %     \begin{itemize}
% %         \item Type of technology and previous experience.
% %         \item Personal risk aversion.
% %     \end{itemize}
% % \end{boxA}

% % \begin{boxB}{Activity VI (Optional): Adopter analysis}
% %     \textbf{Instructions:}\\
% %     Analyze your profile according to the \textbf{Rogers model} based on:
% %     \begin{itemize}
% %         \item Type of technology and previous experience.
% %         \item Personal risk aversion.
% %     \end{itemize}
% % \end{boxB}
